\begin{table}[htbp]
\centering
\small
\caption{Construct validity and reliability of the TDDI.}
\label{tab:validity}
\begin{threeparttable}
\begin{tabular}{lccc}
\toprule
 & Pearson correlation & Spearman correlation & Observations \\
\midrule
\multicolumn{4}{l}{Panel A: Correlations with external proxies and a placebo} \\
Selling-expense ratio & -0.005 & 0.119 & 288 \\
Gross margin & 0.031 & 0.120 & 288 \\
ROA & 0.035 & 0.007 & 289 \\
Firm size & 0.116 & 0.188 & 289 \\
Leverage & -0.001 & -0.067 & 289 \\
Generic digitalization & 0.323 & 0.370 & 298 \\
\midrule
\multicolumn{4}{l}{Panel B: Reliability} \\
Split-half correlation (mean of 200 random splits) & 0.677 & & 298 \\
Spearman--Brown corrected correlation & 0.807 & & 298 \\
\bottomrule
\end{tabular}
\begin{tablenotes}[flushleft]
\footnotesize
\item Notes: External proxies are 2019 firm-level averages. The TDDI is constructed from 2019 annual-report texts. The generic digitalization index is the standardized frequency of category F terms and serves as a placebo. The split-half correlation randomly divides the A--D terms into two halves and correlates the resulting indices across firms.
\end{tablenotes}
\end{threeparttable}
\end{table}
